\begin{proof}[Proof of \cref{obj:lem:variance-envelope-sharpness}]
Write
\[
s=\frac{\sigma_0^2}{n}.
\]
The assumption \(\sigma_0^2\ge0\) gives \(s\ge0\).

\emph{Covariance bound.}
For real random variables \(U,V\) on \(\Omega\) the design covariance is the finite sum
\[
\operatorname{Cov}_\pi(U,V)
=
\sum_{z\in\Omega}\pi(z)\bigl(U(z)-\mathbb E_\pi U\bigr)\bigl(V(z)-\mathbb E_\pi V\bigr),
\]
so applying the Cauchy--Schwarz inequality to the vectors with entries \(\sqrt{\pi(z)}\left(U(z)-\mathbb E_\pi U\right)\) and \(\sqrt{\pi(z)}\left(V(z)-\mathbb E_\pi V\right)\), and then to \(-U\) in place of \(U\), gives
\[
|\operatorname{Cov}_\pi(U,V)|
\le
\sqrt{\operatorname{Var}_\pi(U)}
\sqrt{\operatorname{Var}_\pi(V)}.
\] % lean: variance_envelope_sharpness
Applying the round variance envelope of \cref{obj:ass:round-mean-variance-envelope} to \(X_i\) and \(X_j\), which bounds \(\operatorname{Var}_\pi(X_i)\le s\) and \(\operatorname{Var}_\pi(X_j)\le s\), we obtain
\[
|\operatorname{Cov}_\pi(X_i,X_j)|
\le
\sqrt{\operatorname{Var}_\pi(X_i)}\sqrt{\operatorname{Var}_\pi(X_j)}
\le
\sqrt{s}\sqrt{s}
=
s
\qquad\text{for all }i,j.
\] % lean: variance_envelope_sharpness

\emph{Variance bound.}
By bilinearity of the design covariance, the variance of a linear combination expands as a double sum, and bounding each term by its absolute value gives
\[
\operatorname{Var}_\pi\!\left(\sum_{j=0}^k w_jX_j\right)
=
\sum_{i=0}^k\sum_{j=0}^k
w_iw_j\,\operatorname{Cov}_\pi(X_i,X_j)
\le
\sum_{i=0}^k\sum_{j=0}^k |w_i|\,|w_j|\,s
=
s\left(\sum_{j=0}^k |w_j|\right)^2,
\]
where the inequality uses \(w_iw_j\operatorname{Cov}_\pi(X_i,X_j)\le|w_i||w_j||\operatorname{Cov}_\pi(X_i,X_j)|\le|w_i||w_j|s\) termwise, and the last equality factors the double sum. % lean: variance_envelope_sharpness
This is the stated variance bound.

\emph{Sharpness.}
Define signs
\[
\varepsilon_j=
\begin{cases}
1,& w_j\ge0,\\
-1,& w_j<0,
\end{cases}
\qquad j=0,\ldots,k,
\]
and set
\[
\Gamma_{ij}=s\,\varepsilon_i\varepsilon_j.
\]
This matrix is positive semidefinite: for every \(v\in\mathbb R^{k+1}\),
\[
\sum_{i=0}^k\sum_{j=0}^k v_i\Gamma_{ij}v_j
=
s\left(\sum_{j=0}^k v_j\varepsilon_j\right)^2
\ge0,
\]
since \(s\ge0\). Moreover \(\varepsilon_j^2=1\), so
\[
\Gamma_{jj}=s\,\varepsilon_j^2=s=\frac{\sigma_0^2}{n}
\]
for every \(j\), which in particular satisfies the required diagonal bound. Finally, \(w_j\varepsilon_j=|w_j|\) for every \(j\) by the definition of \(\varepsilon_j\), so the same quadratic-form computation with \(v=w\) gives
\[
\sum_{i=0}^k\sum_{j=0}^k w_i\Gamma_{ij}w_j
=
s\left(\sum_{j=0}^k w_j\varepsilon_j\right)^2
=
s\left(\sum_{j=0}^k |w_j|\right)^2.
\] % lean: variance_envelope_sharpness
Substituting \(s=\sigma_0^2/n\) gives the desired matrix and equality. % lean: variance_envelope_sharpness
\end{proof}
